\documentclass[11pt]{article}
\usepackage[T1]{fontenc}
\usepackage[utf8]{inputenc}
\usepackage{amsmath,amssymb,amsthm}
\usepackage[margin=1in]{geometry}
\newtheorem{theorem}{Theorem}
\title{A Constant Prime Orbit for TLMC 00000000118}
\author{Prepared for the account \texttt{idealistichacker}}
\date{}
\begin{document}
\maketitle

\section*{Frozen statement and explicit reading}
The organizer source at commit
\texttt{95acb520ec5607c826b8a997b1ef2fc82d6f7c57} states in English:
\begin{quote}\raggedright
``For every $M$ there exists an integer-coefficient quadratic $f$\\
such that $f(0),f^1(0),\ldots,f^M(0)$ are all prime.'
\end{quote}
The Chinese source was separately read and gives the same displayed condition.
The frozen Git blob is
\texttt{697131a7b8dc0a62c539018c7d22612df1e7909a}; the SHA-256 of its raw
Git-blob bytes is
\texttt{91b64b4a671db1806c346247acbc2e0d849c5f73337bc807806c7cd00bbd580a}.

We read superscripts as compositional iteration:
\[
 f^{\circ0}(x)=x,\qquad f^{\circ(n+1)}(x)=f\bigl(f^{\circ n}(x)\bigr).
\]
The written list starts with $f(0)$, not $f^{\circ0}(0)=0$. Accordingly, the
formal statement asks for $f(0)$ and for $f^{\circ i}(0)$ when
$1\leq i\leq M$. It does not require the nonprime number $0$ to be prime.
For $M=0$, the resulting requirement is $f(0)$ alone. The source does not
impose pairwise distinctness of the displayed values.

\section*{Construction}
Set
\[
  f(x)=x^2-2x+2=(x-1)^2+1.
\]
This has integer coefficients and nonzero quadratic coefficient $1$. Moreover,
\[
  f(0)=2\quad\text{and}\quad f(2)=2.
\]
Hence every positive iterate at zero is the same prime.

\begin{theorem}
For every natural number $M$, the displayed prime condition holds for the
integer-coefficient quadratic $f(x)=x^2-2x+2$ under the explicit reading above.
\end{theorem}
\begin{proof}
The initial displayed value is $f(0)=2$, which is prime. For each $i\geq1$,
write $i=n+1$. Induction on $n$ shows $f^{\circ(n+1)}(0)=2$: the base case is
$f(0)=2$, and the induction step applies $f(2)=2$. Thus every requested value
is $2$, so it is prime. The same $f$ works for each $M$.
\end{proof}

\section*{Lean correspondence and audit}
The dependency-free Lean 4.33.1 project formalizes the construction in
\texttt{lean4/Main.lean}. Its final theorem is
\begin{center}
\texttt{PrimeChain.iterated\_prime\_chains : forall M : Nat, HasPrimeChain M}.
\end{center}
The predicate \texttt{IntegerQuadratic} explicitly exhibits integer
coefficients and proves the leading coefficient nonzero. \texttt{HasPrimeChain}
encodes the index range stated above. \texttt{lean4/Check.lean} prints the axiom
footprint for every theorem advertised here. The proof does not rely on a
finite computation in place of the quantified theorem.

\section*{Status and boundary}
This PDF was compiled locally with cached-only Tectonic~0.15.0. The package has
no independent review record and makes no claim of organizer approval, merge,
acceptance, rank, award, or human verification. Independent statement-alignment
review, mathematical review, Lean review, reproducibility review, and fresh
release-coordinator checks remain mandatory before any publication.
\end{document}
